\chapter{Kildekode} \label{kap:kode}

\texttt{listings}-pakken er særdeles nyttig, hvis man skal have
kildekode med i sin tekst. Listing \ref{listx} viser et første lille
eksempel på hvad man kan.


\begin{lstlisting}[caption={Et eksempel på Java-kode},label=listx,language=Java]
 public class Main {
  public static void main(String[] args) {
    int time = 16;

    if (time < 12) {
      System.out.println("Dingo.");
    } else if (time < 18) {
      System.out.println("Mango.");
    } else {
      System.out.println("Bingo.");
    }
  }
}
\end{lstlisting}

Man kan også definere sine egne environments, så kildekode i bestemte
sprog bliver vist med reserverede ord fremhævet.

%%% Local Variables:
%%% mode: LaTeX
%%% TeX-master: "report"
%%% End:
